% GENERATED_BY: oc_core_monograph_translation_draft_pipeline_v1.py
% AI_ASSISTED_TRANSLATION_DRAFT: TRUE
% THEOREM_REVIEW_STATUS: PENDING
% THEOREM_REVIEW_MODE: FORMAL_AI_GATE__CODEX_CLI_CHATGPT
% THEOREM_REVIEW_REVIEWER_ID: 
% THEOREM_REVIEW_AUTHORITY_CLASS: 
% THEOREM_REVIEWED_AT: 
% THEOREM_REVIEW_DOSSIER_REF: 
% SOURCE_LANGUAGE: EN
% TARGET_LANGUAGE: DE
% SOURCE_EN_REF: content/19_oc_core_1_3_foundational_consistency.tex
\section{Dossier zur grundlegenden Konsistenz von Core 1.3}
\label{sec:oc13-foundational-consistency}

Dieses Kapitel legt den internen Konsistenzvertrag des Flaggschiff-Bandes fest.
Sein Zweck ist nicht, die formalen Kapitel durch redaktionelle Kommentierung zu
ersetzen, sondern drei Dinge an einer Stelle lesbar zu machen:

\begin{itemize}
\item warum die gegenwärtige Hierarchie als \(K_0\) bis \(K_{12}\) festgelegt
      ist und nicht durch einen beliebigen Schnitt;
\item die vollständige Disziplin der Aussageklassen für axiomatische,
      abgeleitete, empirische, strukturelle, interpretative und Frontier-Aussagen;
\item wie die Master-Monographie und die engere Journal-Core-Extraktion
      zusammenhängen, ohne stillschweigende Drift im Umfang oder in der Notation.
\end{itemize}

\subsection{Warum sich die Hierarchie bis \texorpdfstring{$K_{12}$}{K12} erstreckt}

Core 1.2 enthält bereits genug Struktur, um eine Release-Hierarchie größer als
\(K_{10}\) zu motivieren. Der Punkt ist nicht, dass höhere Bezeichnungen
rhetorisch attraktiv sind. Der Punkt ist, dass der geprüfte Quellkorpus bereits
unterscheidet:

\begin{enumerate}
\item untere Continua, deren Achsen, Schwellen und Flüsse intern zu einer
      gegebenen Modellklasse gehören;
\item theoretische und metatheoretische Continua, die Theorien und
      Modellräume selbst organisieren;
\item eine weitere Schicht, in der Operatoren, Meta-Logiken und Modellräume
      zu expliziten Objekten struktureller Evolution werden;
\item und eine obere Begrenzungsschicht, in der globale Kohärenzbedingungen
      über solche höhergeordneten Strukturen selbst dargestellt werden müssen.
\end{enumerate}

Darum führt Core 1.3 die öffentliche Hierarchie bis \(K_{12}\) und nicht
bei \(K_{10}\) an. Das ist eine Begründung auf Release-Ebene, kein
eigenständiger Unmöglichkeitsbeweis. Der unterstützende Weg ist die
K-Level-Doktrin in Abschnitt~\ref{sec:klevels-full}, die Diskussion des
Theorem-Navigationspfads in Abschnitt~\ref{sec:oc13-theorem-roadmap} und der
durchgespielte Erhaltungstest in Abschnitt~\ref{sec:oc13-worked-examples}.
Diese Abschnitte ordnen \(K_{11}\) den sich entwickelnden metatheoretischen
Objekten und \(K_{12}\) den globalen Kompatibilitäts- und
Integrierbarkeitsbedingungen zu. Die vorliegende Veröffentlichung behauptet
nicht, dass prinzipiell keine weitere Erweiterung denkbar wäre; sie behauptet
nur, dass \(K_{12}\) die höchste Ebene ist, die der gegenwärtige geprüfte
Quellkorpus rechtmäßig trägt, und daher die letzte öffentliche Ebene dieser
Veröffentlichung.

\subsection{Wissenschaftliche Rolle der oberen Ebenen}

\paragraph{Ebene \texorpdfstring{$K_{11}$}{K11}.}
\(K_{11}\) ist die erste Ebene, auf der formale Ontologien, Operatorsysteme,
Meta-Logiken und Modellräume selbst als strukturierte Continua behandelt
werden, die Admissibilitäts-, Transformations- und Kollapsbedingungen
unterliegen.
Diese Ebene wird nötig, sobald man nicht nur Theorien untersucht, sondern die
Entwicklung der Räume, in denen Theorien und theorieerzeugende Operatoren leben.

\paragraph{Ebene \texorpdfstring{$K_{12}$}{K12}.}
\(K_{12}\) wird nicht als mystische Vollendungsschicht eingeführt.
Es ist die minimal spezifizierte obere Kohärenzschicht, die benötigt wird,
wenn man globale Kompatibilität, Beschränktheit und Fehlermodi über Familien
von \(K_{11}\)-Objekten hinweg diskutieren will.
Seine Rolle ist strukturell und nicht triumphalistisch: Es verhindert, dass
die Hierarchie so tut, als gäbe es oberhalb rekursiv evolvierender
metatheoretischer Räume kein formales Problem mehr.

\subsection{Aussageklassen}

Um die Monographie lesbar zu halten, ohne den wissenschaftlichen Status zu
verwischen, verwendet Core~1.3 in diesem Kapitel das folgende
Beweisstatus-Vokabular. Es ist die lokale Taxonomie von Theorem- und
Anspruchsumfang; Abschnitt~\ref{sec:oc13-theorem-roadmap} ergänzt die
nachgelagerte operative Inspektionslast für Replay-, Comparator-, Falsifikator-
und Eskalationsregeln.

\begin{itemize}
\item \textbf{Axiomatische Aussagen} fixieren primitive Annahmen,
      Admissibilitätsbedingungen oder Identitätsregeln.
\item \textbf{Abgeleitete Aussagen} sind theoremtragende Aussagen mit expliziter
      Unterstützungsabhängigkeit und Beweispflichten.
\item \textbf{Empirische Aussagen} verlangen ein eigenständiges,
      quellenbesitzendes Domänenpaket, eine kanonische theorem-native Spur,
      eine gepinnte Datenroute, eine numerische Parameterregel, eine
      Observable-Bindung, ein gesperrtes Held-out-Replay-Verfahren, ein
      Comparator-Audit und einen expliziten Falsifikator vor der Promotion.
\item \textbf{Strukturelle Aussagen} beschreiben die Organisation des Rahmens,
      die Hierarchie und die Abbildung zwischen Modulen oder Ebenen.
\item \textbf{Interpretative Aussagen} erklären, was die formale Struktur
      bedeutet, für welchen Leser sie gedacht ist und was spätere Projektionen
      testen dürfen.
\item \textbf{Frontier-Aussagen} dürfen im Forschungsprogramm verbleiben,
      dürfen aber ohne formale Promotion nicht in positive äußere Wissenschaft
      eingehen.
\end{itemize}
Diese lokalen Klassen sind nicht als disjunkt vorausgesetzt. Wenn ein Satz
mehr als einer Klasse entsprechen könnte, wird sein öffentlicher Status nach
semantischer Rolle und Beweisfunktion zugewiesen, nicht nach der
Oberflächenformulierung:
\begin{enumerate}
\item ein expliziter Widerspruch zu einem promoteten Anspruch markiert
      zunächst einen Refutationskandidaten oder einen Oberflächenwiderspruch
      für das SPOT-Audit;
\item empirischer Gehalt wird vor jedem Vertrauen in die Prosa anhand von
      Paket, Datenroute, Replay, Comparator und Falsifikator geprüft;
\item abgeleiteter Theorem-Gehalt wird entlang seiner Beweisroute und seines
      Abhängigkeitsgraphen geprüft;
\item zugelassene primitive Annahmen bleiben axiomatische Prämissen;
\item strukturelle Abbildungs- und Hierarchieaussagen bleiben strukturell,
      auch wenn sie eine mögliche spätere Nutzung erwähnen;
\item erklärende, leserorientierte oder didaktische Sätze bleiben interpretativ,
      sofern sie nicht theorematische oder empirische Kraft behaupten;
\item nicht promotete Forschungsprogramm-Vorschläge bleiben Frontier-Aussagen;
\item rein struktureller oder interpretativer Gehalt behält die von SPOT
      protokollierte Klasse \occode{FRAME_ONLY} oder \occode{EXTENSION},
      sofern er nicht explizit an evidenztragendes Bridge-Material gekoppelt ist;
\item \occode{BRIDGE_ONLY} ist eine stabile evidenzstützende Klasse für Zeilen
      oder Pakete, die Brückenstütze unter expliziten Umfangsnotizen tragen;
      sie ist keine Warteschlange für theorem-native Promotion.
\end{enumerate}
Wörter wie ``frontier'', ``Kandidat'' oder ``Projektion'' sind daher
diagnostische Signale, nicht die eigentliche Klassenregel.

Diese Prosa-Klassen sind eine lokale Beweisstatus-Taxonomie, kein Ersatz für
das kanonische SPOT-Vokabular der wissenschaftlichen Klassen und
Abschlussverdicts. Release-Verdikte wie \texttt{PASS} und
\texttt{FAIL\_CLOSED} bleiben eine separate Verdict-Schicht. Die Zuordnung ist:
\begin{itemize}
\item Axiomatische Aussagen sind zugelassene Annahmen oder Identitätsregeln.
      Sie können als benannte Annahmen in einer \occode{THEOREM_NATIVE}-Route
      stehen, aber das Axiom selbst wird dadurch nicht zu einem bewiesenen
      Theorem.
\item Abgeleitete Aussagen werden auf \occode{THEOREM_NATIVE} abgebildet,
      wenn die Unterstützungsroute geschlossen ist.
\item Empirische Aussagen werden nur nach Durchlaufen von quellenbesitzendem
      Paket, theorem-native Spur, gepinnter Datenroute, numerischer
      Parameterregel, Observable-Bindung, gesperrtem Held-out-Replay-Verfahren,
      Comparator-Audit und Falsifikator-Gates auf \occode{THEOREM_NATIVE}
      abgebildet.
\item Strukturelle und interpretative Aussagen werden auf \occode{FRAME_ONLY}
      oder \occode{EXTENSION} abgebildet, wie von SPOT protokolliert, sofern
      die Aussage nicht an eine evidenztragende Bridge-Zeile gekoppelt ist.
\item Frontier-Aussagen werden auf \occode{EXTENSION} abgebildet, bis ein
      späterer Abschnitt oder eine spätere Veröffentlichung sie durch dieselben
      theorem-native oder empirischen Gates promotet, die auch für andere
      positive Aussagen verwendet werden, oder eine explizite Refutation
      dokumentiert.
\item Fail-closed Bridge-, Frame- oder Extension-Lanes behalten die von SPOT
      protokollierte Klasse; ein Closure-Fehlschlag ist für sich genommen keine
      Refutation.
\item Aussagen werden nur dann auf \occode{REFUTED} abgebildet, wenn SPOT eine
      explizite Refutation statt einer blockierten Promotion protokolliert.
\end{itemize}
Dieses Kapitel klassifiziert die Aussageklassen, die Theorem-Status,
strukturellen Umfang, Interpretation, empirische Promotion und Frontier-Grenzen
beeinflussen. Die operative Routendisziplin wird später in den Kapiteln zur
Theorem-Roadmap und zur Beweismaschine behandelt.
Sprache der redaktionellen Brücke ist keine siebte Beweisklasse; sie ist eine
interpretative und release-konsistente Kennzeichnung für Aussagen, die
Monographie, Journal Core, Dossiers und das öffentliche Paket verbinden, ohne
Theoremkraft hinzuzufügen.
Operative Aussagen sind die späteren exakten Abgleich-, Held-out-Replay-,
statistischen Schwellen-, Tail-Breach-, Comparator-, Falsifikator- und
Eskalationsregeln, die an die empirische Promotion gebunden sind.
Strukturelle Aussagen werden stattdessen anhand von Umfangs-, Abhängigkeits-
und Quellenrandprüfungen beurteilt. Beide Kontrollarten werden erst dann
behandelt, nachdem hier bereits die Aussageklasse, die numerische
Parameterregel und die Observable-Bindung festgelegt worden sind; keine von
beiden wird zu einer siebten Beweisstatusklasse.

Die Journal-Core-Extraktion darf diesen Bestand verengen, aber nicht diese
Klassen ineinander verwischen.
Ein interpretativer Satz kann ein Theorem erläutern; er darf das Theorem nicht
ersetzen.
Eine strukturelle Bemerkung kann eine Projektion motivieren; sie darf nicht
stillschweigend als Beweis dieser Projektion dienen.

\subsection{Beweistragende und Nichtbehauptungsschichten}

Die Master-Monographie trennt daher drei Schichten explizit:

\begin{enumerate}
\item das übernommene formale Rückgrat, das noch immer die Struktur der
      Kerntheorie trägt;
\item die Reparaturschicht von Core 1.3, die Zeugendisziplin, Theorem-Schicksal,
      Chronologie und Veröffentlichungsgrenze klärt;
\item die begrenzte Nichtbehauptungsschicht, die festhält, was die vorliegende
      Veröffentlichung noch nicht verteidigt, auch wenn das breitere
      Forschungsprogramm es weiterhin verfolgt.
\end{enumerate}

Diese Trennung ist für feindselige Begutachtung wichtig.
Ein Gutachter sollte für jeden Satz, der stark klingt, fragen können:
Ist er axiomatisch, abgeleitet, empirisch, strukturell, interpretativ oder
Frontier?
Wenn er abgeleitet ist, was stützt ihn?
Wenn er empirisch ist, welches Paket, welche Datenroute, welches Replay,
welcher Comparator und welcher Falsifikator stützen ihn?
Wenn er strukturell oder interpretativ ist, was behauptet er genau nicht?
Wenn er Frontier ist, wo ist sichtbar verankert, dass er keine positive
Veröffentlichungsautorität hat?

\subsection{Beziehung zwischen Master-Monographie und Journal Core}

Der kanonische SPOT ist die wissenschaftliche Autorität der Veröffentlichung,
und der quellenbesitzende Wissenschaftskorpus ist der Autorierungs- und
Zeugenort hinter dieser Projektion. Die Master-Monographie ist die
Flaggschiff-Referenz für Leser, weil sie die vollständige Theorie, die
Leserleiter, die Anhänge und die explizite Veröffentlichungsgrenze an einer
Stelle trägt.
Der \emph{Journal Core} ist das begrenzte Artikelartefakt. Die
\emph{Journal-Core-Extraktion} ist die Operation, die einen defensiblen
theoremtragenden Pfad zur redaktionellen Prüfung auswählt, ohne so zu tun, als
sei die gesamte Monographie bereits in dieses kürzere Artefakt komprimiert.

Aus diesem Grund ist die Extraktionsregel nur in eine Richtung gültig:

\[
\text{Master Monographie} \longrightarrow \text{Journal Core},
\]

mit expliziter Erhaltung von Notation, Nichtbehauptungsgrenzen und Theoremumfang.
Die Extraktion darf Details weglassen. Der Journal Core darf den Anspruch nicht
vergrößern.
Die öffentliche Abschlussgrenze bleibt der in
Abschnitt~\ref{sec:oc13-publication-boundary} genannte Umfang
\texttt{CORE\_1\_3\_SCIENCE\_ONLY}. \texttt{FRAME\_ONLY}-Abschlüsse dürfen als
Frame-Stütze \texttt{PASS} sein, werden dadurch aber nicht zu eigenständigen
empirischen Promotions. \texttt{BRIDGE\_ONLY}-Zeilen sind an den Umfang
gebundene evidenzielle Stütze und keine öffentlich promotete Abschlussklasse;
sie werden nicht als theorem-native \texttt{PASS}-Zeilen neu etikettiert, es
sei denn, es wird ein separates quellenbesitzendes theorem-native Paket
angelegt und besteht die kanonischen Gates. \texttt{EXTENSION}- und
\texttt{REFUTED}-Lanes bleiben außerhalb der öffentlichen \texttt{PASS}-Autorität.

\subsection{Konsequenz für Leser}

Dieses Kapitel existiert, weil wissenschaftliche Strenge und Lesbarkeit hier
nicht gegeneinander arbeiten.
Je expliziter das Dossier über Hierarchie, Aussageklasse und Beweisgrenze wird,
desto leichter wird es gleichzeitig für drei Leser:

\begin{itemize}
\item den Editor, der einen sauberen Überblick darüber braucht, was tatsächlich
      verteidigt wird;
\item den formalen Leser, der exakten Umfang und strenge Abhängigkeitsdisziplin
      benötigt;
\item den nicht spezialisierten wissenschaftlichen Leser, der einen kohärenten
      narrativen Einstiegspunkt braucht, bevor er in die vollständige technische
      Hülle eintaucht.
\end{itemize}

\paragraph{Brücke zum nächsten Kapitel.}
Mit den festgelegten Aussageklassen, der Hierarchiegrenze und der Beziehung von
Master zu Derivat beschränkt sich das nächste Kapitel auf genau eine Aufgabe:
den kürzesten zulässigen Weg der Theorem-Navigation durch das tragende
Theorem-Rückgrat.

\clearpage
